% !TeX program = xelatex
\documentclass[10pt,aspectratio=169]{beamer}

\usepackage{ctex}
\usepackage{metalogo}
\usepackage{hznu}

\usefonttheme[onlymath]{serif}

\title[杭州师范大学 Beamer 模板]{杭州师范大学 Beamer 模板}
\subtitle{——这里是副标题}

\author[Yeung Zhaam]{%
  \texorpdfstring{%
    杨湛 \\\medskip
    {\small
      \href{mailto:2025212501117@stu.hznu.edu.cn}
        {\nolinkurl{<2025212501117@stu.hznu.edu.cn>}}%
    } \\
    {\small \url{https://www.hznu.edu.cn/}}%
  }{杨湛}%
}

\institute[数学学院]{%
  数学学院 \\
  杭州师范大学%
}

\date[Apr. 4, 2026]{2026年4月4日}

\begin{document}

\begin{frame}[plain]
  \titlepage
\end{frame}

\section{提纲}
\begin{frame}{提纲}
  \tableofcontents
\end{frame}

\section{介绍}
\begin{frame}{介绍}
  \begin{itemize}
    \item 编译方式
      \begin{itemize}
        \item 推荐安装完整版的 TeX Live
        \item 使用 \XeLaTeX 编译
      \end{itemize}
    \item 请参考 \LaTeX 和 Beamer 用户文档
    \item 行内数学公式示例 $\sin^2 \theta + \cos^2 \theta = 1$
    \item 行间数学公式示例
      \begin{equation}
        y_1 = \int \sin x\, \mathrm{d}x
      \end{equation}
    \item 基于“师大蓝”颜色：\url{https://www.hznu.edu.cn/}
  \end{itemize}
\end{frame}

\section{内置环境}
\begin{frame}{内置环境}
  \begin{block}{Slides with \LaTeX}
    Beamer offers a lot of functions to create nice slides using \LaTeX.
  \end{block}

  \begin{block}{The basis}
    内部使用以下主题：
    \begin{itemize}
      \item split
      \item whale
      \item rounded
      \item orchid
    \end{itemize}
  \end{block}
\end{frame}

\begin{frame}{带数字列表}
  \begin{enumerate}
    \item This just shows the effect of the style
    \item It is not a Beamer tutorial
    \item Read the Beamer manual for more help
    \item Contact me only concerning the style file
  \end{enumerate}
\end{frame}

\section{结论}
\begin{frame}{结论}
  \begin{itemize}
    \item Easy to use
    \item Good results
  \end{itemize}
\end{frame}

\section{参考文献}
\begin{frame}{参考文献}
  \begin{thebibliography}{99}
    \bibitem{zhao1}
      Yi~Zhao, \textit{An introduction to X}, Sep.~15, 2015
    \bibitem{qian2}
      Er~Qian, San~Sun,
      Phys.\ Lett.\ A \textbf{xx}, 2xx (20xx)
    \bibitem{li4}
      Si~Li, Phys.\ Rev.\ C \textbf{xx}, 5xx (20xx)
  \end{thebibliography}
\end{frame}

\end{document}
